\documentclass[12pt]{article}
\usepackage{graphicx}
\usepackage{amsmath}

\title{Selective Adsorption of Nanoplastic Polymer Fragments in Zeolite MFI}
\author{Author et al.}

\begin{document}
\maketitle

\begin{abstract}
Nanoplastic contamination poses an emerging environmental challenge. Here we investigate the
selective adsorption of polyethylene (PE) and polypropylene (PP) fragments in zeolite MFI using
density functional theory (DFT) and grand canonical Monte Carlo (GCMC) simulations. PE fragments
adsorb preferentially over PP, with a selectivity of approximately 2.0 under ambient conditions.
\end{abstract}

\section{Introduction}
Nanoplastics released from weathering of plastic debris accumulate in water systems.
Zeolites offer potential as low-cost adsorbents due to their tunable pore architecture.
Zeolite MFI has pore dimensions (5.3 × 5.6 Å) that are comparable to small polymer
fragment chain widths, motivating a molecular-scale investigation of adsorption selectivity.

\section{Methods}

\subsection{DFT calculations}
All periodic DFT calculations were performed using CASTEP. Geometry optimisations used the
generalised gradient approximation with dispersion corrections. A plane-wave energy cutoff was
applied. Basis set superposition error (BSSE) was corrected using the counterpoise method.

\subsection{Monte Carlo simulations}
Grand canonical Monte Carlo (GCMC) simulations were performed using RASPA2. PE and PP
trimers were modelled using a united-atom force field. Simulations were run at 298 K
over a range of pressures from 0.01 to 1 bar.

\subsection{Pore characterisation}
Geometric pore characteristics including largest included sphere diameter and pore-limiting
diameter were calculated using Zeo++.

\section{Results}

\subsection{Adsorption energies}
The most stable adsorption site for PE trimers is the straight-channel intersection.
BSSE-corrected adsorption energies range from −42 to −61 kJ/mol across five investigated sites.
PP adsorption energies are systematically lower by approximately 26 kJ/mol.

\subsection{Isotherms}
At 0.1 bar and 298 K, PE loading reaches approximately 4.2 molecules per unit cell,
compared to 2.1 for PP. The resulting selectivity of $S_{\rm PE/PP} \approx 2.0$ persists
across the full pressure range investigated.

\section{Discussion}
The preferential adsorption of PE over PP can be attributed to the better geometric
compatibility of the linear PE chain with the MFI straight-channel geometry. Dispersion
interactions, which scale with contact area, favour the planar PE backbone over the
branched PP structure.

\section{Conclusions}
MFI zeolite shows a 2-fold selectivity for PE over PP nanoplastic fragments under
ambient conditions. This result suggests potential application of zeolites as
size-selective nanoplastic scavengers.

\end{document}
